\documentclass[a4paper,11pt]{article}

\usepackage[T1]{fontenc}
\usepackage[utf8]{inputenc}
\usepackage[slovene]{babel}
\usepackage[margin=1.5cm]{geometry}
\usepackage{tabularx}
\usepackage{array}

\title{Moj tedenski urnik}
\author{Gaber Košir}
\date{9.--15. november 2026}

\begin{document}
\maketitle

\section*{Urnik predavanj in vaj -- FMF}

\small
\renewcommand{\arraystretch}{1.5}

\begin{tabularx}{\textwidth}{|l|l|X|}
\hline
\textbf{Dan} & \textbf{Ura} & \textbf{Predmet} \\
\hline
Ponedeljek & 9--12 & Linearna algebra (P) \\
& 12--14 & Matematika I (P) \\
& 14--16 & Fizika I (P) \\
\hline
Torek & 9--11 & Linearna algebra (V) \\
& 11--13 & Fizika I (V) \\
& 13--15 & Tutorstvo APF 1. letnik \\
\hline
Sreda & 11--12 & Linearna algebra (V) \\
& 12--14 & Matematika I (V) \\
& 14--17 & Praktikum aplikativne fizike I \\
\hline
Četrtek & 8--10 & Fizika I (P) \\
& 10--11 & Fizika I (S) \\
& 11--12 & Fizika I (V) \\
& 12--14 & Matematika I (P) \\
& 15--16 & Konstrukcijski elementi (P) \\
& 16--17 & Konstrukcijski elementi (V) \\
\hline
Petek & 9--11 & Matematika I (V) \\
& 11--12 & Računalniški praktikum (P) \\
& 12--15 & Računalniški praktikum (V) \\
\hline
\end{tabularx}

\vspace{1cm}

\noindent
P = predavanja, V = vaje, S = seminar.

\end{document}
